\lettersection{Academic Experience} % your preparation for the field of study

I've learned how to program in C++ (CS902 98) and taught myself with other programming languages such as Python and JavaScript.
I used what I've learned to solve practical problems.
I built some tiny tools to download the class table in iCalendar format, to classify PDF pages by page size, or to monitor an SMS platform for replies.
I also participated in a web project which aims to provide an office automation system for students' network managers.
During these projects, I improved my programming skills and learned how to build applications based on the need of users and how to collaborate with others.

In the field of security, I've also learned basic topics in security (IS304 97, IS497 96, NUS CS3235 A).
And I partially got some knowledge of system security (IS405 92) and formal perspective of cryptography (NUS CS4236 A+).
In the National College Student Information Security Contest, our group developed a blockchain-based online shopping platform oriented to protecting user privacy, which won the second prize.

I act as a research intern in \textit{\textbf{Lab of Cryptography and Computer Security}} in SJTU in September 2018, where I got my first experience with cryptography theory, network essentials and basic concepts of Bitcoin and Blockchain. By giving and listening to reports, I learned how to organize a presentation.
Under the guidance of the professor, I completed a project independently to integrate digest algorithm SM3 to a modified version of OpenSSL. This project led me to a deeper exploration of OpenSSL.
I also helped build a proof-of-concept project, in which we improved consensus efficiency by sharding and two-layer consensus.

During my internship in \textit{\textbf{National Cybersecurity R\&D Lab}} in the fall of 2019, where I was supervised by Chang Ee-Chien, I learned Intel Software Extension (SGX).
In that semester, I proposed a method of mutual attestation between enclaves which significantly reduced the overhead on communication.
And I successfully migrated an open-source implementation of Raft into SGX, which utilized the mutual attestation protocol in a multi-node network.

In this semester, I'm doing my Final Year Project continuing my interest in Blockchain and SGX.
I am required to design a new smart contract system supporting secure multiparty computation utilizing the consistency provided by blockchain and security guarantee provided by SGX.

I believe I can complete the programme well, both in courses and in research.